\section*{Exercises on \S\arabic{exercisegroup}}
\phantomsection
\addcontentsline{toc}{section}{Exercises on \protect\S\arabic{exercisegroup}}

\begin{enumerate}
\def\labelenumi{\arabic{enumi})}
\tightlist
\item
  With the same hypotheses as in prop. 2 of I, p.~20 prove the formula
\end{enumerate}

\[
[ \mathbf {f} ^ {(n)}. \mathbf {g} ] = \sum_ {p = 0} ^ {n} (- 1) ^ {p} \binom {n} {p} D ^ {n - p} [ \mathbf {f}. \mathbf {g} ^ {(p)} ].
\]

\begin{enumerate}
\def\labelenumi{\arabic{enumi})}
\setcounter{enumi}{1}
\tightlist
\item
  With the notation of prop. 2 of I, p.~28 suppose that the relation \([\mathbf{a}.\mathbf{y}] = 0\) for all \(\mathbf{y} \in \mathrm{F}\) implies that \(\mathbf{a} = 0\) in E. Under these conditions, if \(\mathbf{g}_i\) ( \(0 \leqslant i \leqslant n\) ) are \(n + 1\) vector functions with values in E, defined on an interval I of R and such that for every vector function f with values in F and n times differentiable on I, one has identically
\end{enumerate}

\[
[ \mathbf {g} _ {0}. \mathbf {f} ] + [ \mathbf {g} _ {1}. \mathbf {f} ^ {\prime} ] + \dots + [ \mathbf {g} _ {n}. \mathbf {f} ^ {(n)} ] = 0
\]

then the functions \(g_{i}\) are identically zero.

\begin{enumerate}
\def\labelenumi{\arabic{enumi})}
\setcounter{enumi}{2}
\tightlist
\item
  With the notation of exerc. 2 and the same hypothesis on \([x.y]\) suppose that each of the functions \(g_{k}\) is n times differentiable on I; for each function f which is n times differentiable on I, with values in F, put
\end{enumerate}

\[
[ \mathbf {g} _ {0}. \mathbf {f} ] - [ \mathbf {g} _ {1}. \mathbf {f} ] ^ {\prime} + [ \mathbf {g} _ {2}. \mathbf {f} ] ^ {\prime \prime} + \dots + (- 1) ^ {n} [ \mathbf {g} _ {n}. \mathbf {f} ] ^ {(n)} = [ \mathbf {h} _ {0}. \mathbf {f} ] + [ \mathbf {h} _ {1}. \mathbf {f} ^ {\prime} ] + \dots + [ \mathbf {h} _ {n}. \mathbf {f} ^ {(n)} ],
\]

which defines the functions \(\mathbf{h}_i\) ( \(0 \leqslant i \leqslant n\) ) without ambiguity (exerc. 2); show that one has

\[
[ \mathbf {h} _ {0}. \mathbf {f} ] - [ \mathbf {h}. \mathbf {f} ] ^ {\prime} + [ \mathbf {h} _ {2}. \mathbf {f} ] ^ {\prime \prime} + \dots + (- 1) ^ {n} [ \mathbf {h} _ {n}. \mathbf {f} ] ^ {(n)} = [ \mathbf {g} _ {0}. \mathbf {f} ] + [ \mathbf {g} _ {1}. \mathbf {f} ^ {\prime} ] + \dots + [ \mathbf {g} _ {n}. \mathbf {f} ^ {(n)} ]
\]

identically.

\begin{enumerate}
\def\labelenumi{\arabic{enumi})}
\setcounter{enumi}{3}
\tightlist
\item
  Let \(\mathbf{f}\) be a vector function which is \(n\) times differentiable on an interval \(\mathrm{I} \subset \mathbf{R}\) . Show that for \(1 / x \in \mathrm{I}\) one has identically
\end{enumerate}

\[
\frac {1}{x ^ {n + 1}} \mathbf {f} ^ {(n)} \left(\frac {1}{x}\right) = (- 1) ^ {n} \mathrm{D} ^ {n} \left(x ^ {n - 1} \mathbf {f} \left(\frac {1}{x}\right)\right)
\]

(argue inductively on \(n\) ).

\begin{enumerate}
\def\labelenumi{\arabic{enumi})}
\setcounter{enumi}{4}
\tightlist
\item
  Let \(u\) and \(v\) be two real functions which are \(n\) times differentiable on an interval \(I \subset \mathbf{R}\) . If one puts \(D^n(u / v) = (-1)^n w_n / v^{n+1}\) at every point where \(v(x) \neq 0\) , show that
\end{enumerate}

\[
w _ {n} = \left| \begin{array}{c c c c c c} u & v & 0 & 0 & \dots & 0 \\ u ^ {\prime} & v ^ {\prime} & v & 0 & \dots & 0 \\ u ^ {\prime \prime} & v ^ {\prime \prime} & 2 v ^ {\prime} & v & \dots & 0 \\ \dots & \dots & \dots & \dots & \dots & \dots \\ u ^ {(n - 1)} & v ^ {(n - 1)} & \binom {n - 1} {1} v ^ {(n - 2)} & \binom {n - 1} {2} v ^ {(n - 3)} & \dots & v \\ u ^ {(n)} & v ^ {(n)} & \binom {n} {1} v ^ {(n - 1)} & \binom {n} {2} v ^ {(n - 2)} & \dots & \binom {n} {n - 1} v ^ {\prime} \end{array} \right|
\]

(put \(w = u / v\) and differentiate \(n\) times the relation \(u = wv\) ).

\begin{enumerate}
\def\labelenumi{\arabic{enumi})}
\setcounter{enumi}{5}
\tightlist
\item
  Let \(\mathbf{f}\) be a vector function defined on an open interval \(\mathrm{I} \subset \mathbf{R}\) , taking values in a normed space \(\mathrm{E}\) .
\end{enumerate}

Put \(\Delta \mathbf{f}(x; h_1) = \mathbf{f}(x + h_1) - \mathbf{f}(x)\) , and then, inductively, define

\[
\Delta^ {p} \mathbf {f} (x; h _ {1}, h _ {2}, \dots , h _ {p - 1}, h _ {p}) = \Delta^ {p - 1} \mathbf {f} (x + h _ {p}; h _ {1}, \dots , h _ {p - 1}) - \Delta^ {p - 1} \mathbf {f} (x; h _ {1}, \dots , h _ {p - 1});
\]

these functions are defined for each \(x \in I\) when the \(h_{i}\) are small enough.

\begin{enumerate}
\def\labelenumi{\alph{enumi})}
\tightlist
\item
  If the function \(\mathbf{f}\) is \(n\) times differentiable at the point \(x\) (and so \(n - 1\) times differentiable on a neighbourhood of \(x\) ), one has
\end{enumerate}

\[
\lim_{\substack{(h_{1},\ldots ,h_{n})\to (0,\ldots ,0)\\ h_{1}h_{2}\ldots h_{n}\neq 0}}\frac{\varDelta^{n}\mathbf{f}(x;h_{1},\ldots,h_{n})}{h_{1}h_{2}\ldots h_{n}} = \mathbf{f}^{(n)}(x)
\]

(argue by induction on n, using the mean value theorem).

\begin{enumerate}
\def\labelenumi{\alph{enumi})}
\setcounter{enumi}{1}
\tightlist
\item
  If \(\mathbf{f}\) is \(n\) times differentiable on the interval I, one has
\end{enumerate}

\[
\begin{array}{l} \left\| \Delta^ {n} \mathbf {f} (x; h _ {1}, \ldots , h _ {n}) - \mathbf {f} ^ {(n)} (x _ {0}) h _ {1} h _ {2} \ldots h _ {n} \right\| \\ \leqslant | h _ {1} h _ {2} \ldots h _ {n} | \sup \left\| \mathbf {f} ^ {(n)} (x + t _ {1} h _ {1} + \dots + t _ {n} h _ {n}) - \mathbf {f} ^ {(n)} (x _ {0}) \right\| \end{array}
\]

the supremum being taken over the set of \((t_i)\) such that \(0 \leqslant t_i \leqslant 1\) for \(1 \leqslant i \leqslant n\) (same method).

\begin{enumerate}
\def\labelenumi{\alph{enumi})}
\setcounter{enumi}{2}
\tightlist
\item
  If \(f\) is a real function which is \(n\) times differentiable on I, one has
\end{enumerate}

\[
\Delta^ {n} f (x; h _ {1}, h _ {2}, \dots , h _ {n}) = h _ {1} h _ {2} \dots h _ {n} f ^ {(n)} (x + \theta_ {1} h _ {1} + \dots + \theta_ {n} h _ {n})
\]

the numbers \(\theta_{i}\) belonging to {[}0, 1{]} (same method, using I, p.~22, corollary).

\begin{enumerate}
\def\labelenumi{\arabic{enumi})}
\setcounter{enumi}{6}
\tightlist
\item
  Let \(f\) be a finite real function \(n\) times differentiable at the point \(x_0\) , and \(\mathbf{g}\) a vector function which is \(n\) times differentiable at the point \(y_0 = f(x_0)\) . Let
\end{enumerate}

\[
\begin{array}{l} f (x _ {0} + h) = a _ {0} + a _ {1} h + \dots + a _ {n} h ^ {n} + r _ {n} (h) \\ \mathbf {g} (y _ {0} + k) = \mathbf {b} _ {0} + \mathbf {b} _ {1} k + \dots + \mathbf {b} _ {n} k ^ {n} + \mathbf {s} _ {n} (k) \end{array}
\]

be the Taylor expansions of order n of f and g at the points \(x_{0}\) and \(y_{0}\) respectively. Show that the sum of the \(n+1\) terms of the Taylor expansion of order n of the composite function \(g \circ f\) at the point \(x_{0}\) is equal to the sum of the terms of degree \(\leqslant n\) in the polynomial

\[
\mathbf {b} _ {0} + \mathbf {b} _ {1} (a _ {1} h + \dots + a _ {n} h ^ {n}) + \mathbf {b} _ {2} (a _ {1} h + \dots + a _ {n} h ^ {n}) ^ {2} + \dots + \mathbf {b} _ {n} (a _ {1} h + \dots + a _ {n} h ^ {n}) ^ {n}.
\]

Deduce the two following formulae:

\begin{enumerate}
\def\labelenumi{\alph{enumi})}
\tightlist
\item
\end{enumerate}

\[
\mathrm{D} ^ {n} (\mathbf {g} (f (x))) = \sum \frac {n !}{m _ {1} ! m _ {2} ! \dots m _ {q} !} \mathbf {g} ^ {(p)} (f (x)) \left(\frac {f ^ {\prime} (x)}{1 !}\right) ^ {m _ {1}} \dots \left(\frac {f ^ {(q)} (x)}{q !}\right) ^ {m _ {q}}
\]

the sum being taken over all systems of positive integers \((m_{i})_{1\leqslant i\leqslant q}\) such that

\[
m _ {1} + 2 m _ {2} + \dots + q m _ {q} = n
\]

where p denotes the sum \(m_{1} + m_{2} + \cdots + m_{q}\) .

\begin{enumerate}
\def\labelenumi{\alph{enumi})}
\setcounter{enumi}{1}
\tightlist
\item
\end{enumerate}

\[
\mathrm{D} ^ {n} (\mathbf {g} (f (x))) = \sum_ {p = 1} ^ {n} \frac {1}{p !}   \mathbf {g} ^ {(p)} (f (x)) \left(\sum_ {q = 1} ^ {p} \binom {p} {q} (- f (x)) ^ {p - q}   \mathrm{D} ^ {n} ((f (x)) ^ {q}\right).
\]

\begin{enumerate}
\def\labelenumi{\arabic{enumi})}
\setcounter{enumi}{7}
\tightlist
\item
  Let \(f\) be a real function defined and \(n\) times differentiable on an interval \(\mathrm{I}\) , let \(x_{1}, x_{2}, \ldots, x_{p}\) be distinct points of \(\mathrm{I}\) , and \(n_{i}\) ( \(1 \leqslant i \leqslant p\) ) be \(p\) integers \(> 0\) such that
\end{enumerate}

\[
n _ {1} + n _ {2} + \dots + n _ {p} = n.
\]

Suppose that at the point \(x_{i}\) the function f vanishes together with its first \(n_{i}-1\) derivatives for \(1 \leqslant i \leqslant p\) : show that there is a point \(\xi\) interior to the smallest interval that contains the \(x_{i}\) and such that \(f^{(n-1)}(\xi) = 0\) .

\begin{enumerate}
\def\labelenumi{\arabic{enumi})}
\setcounter{enumi}{8}
\tightlist
\item
  With the same notation as in exerc. 8 suppose that \(f\) is \(n\) times differentiable on I but otherwise arbitrary. Let \(g\) be the polynomial of degree \(n - 1\) (with real coefficients) such that at the point \(x_{i}\) ( \(1 \leqslant i \leqslant p\) ) both \(g\) and its first \(n_{i} - 1\) derivatives are respectively equal to \(f\) and its first \(n_{i} - 1\) derivatives. Show that we have
\end{enumerate}

\[
f (x) = g (x) + \frac {(x - x _ {1}) ^ {n _ {1}} (x - x _ {2}) ^ {n _ {2}} \dots (x - x _ {p}) ^ {n _ {p}}}{n !} f ^ {(n)} (\xi)
\]

where \(\xi\) is interior to the smallest interval containing the points \(x_{i}\) ( \(1 \leqslant i \leqslant p\) ) and \(x\) . (Apply exerc. 8 to the function of \(t\)

\[
f (t) - g (t) - a \frac {(t - x _ {1}) ^ {n _ {1}} (t - x _ {2}) ^ {n _ {2}} \dots (t - x _ {p}) ^ {n _ {p}}}{n !}
\]

where \(a\) is a suitably chosen constant.)

\begin{enumerate}
\def\labelenumi{\arabic{enumi})}
\setcounter{enumi}{9}
\tightlist
\item
  Let \(g\) be an odd real function defined on a neighbourhood of 0, and 5 times differentiable on this neighbourhood. Show that
\end{enumerate}

\[
g (x) = \frac {x}{3} \left(g ^ {\prime} (x) + 2 g ^ {\prime} (0)\right) - \frac {x ^ {5}}{1 8 0} g ^ {(5)} (\xi) \quad (\xi = \theta x, \quad 0 <   \theta <   1)
\]

(same method as in exerc. 9).

Deduce that if \(f\) is a real function defined on \([a, b]\) and 5 times differentiable on this interval, then

\[
f (b) - f (a) = \frac {b - a}{6} \left[ f ^ {\prime} (a) + f ^ {\prime} (b) + 4 f ^ {\prime} \left(\frac {a + b}{2}\right) \right] - \frac {(b - a) ^ {5}}{2 8 8 0} f ^ {(5)} (\xi)
\]

with \(a < \xi < b\) (``Simpson's formula'').

\begin{enumerate}
\def\labelenumi{\arabic{enumi})}
\setcounter{enumi}{10}
\tightlist
\item
  Let \(f_{1}, f_{2}, \ldots, f_{n}\) and \(g_{1}, g_{2}, \ldots, g_{n}\) be \(2n\) real functions which are \(n - 1\) times differentiable on an interval I. Let \((x_{i})_{1 \leqslant i \leqslant n}\) be a strictly increasing sequence of points in I. Show that the ratio of the two determinants
\end{enumerate}

\[
\left| \begin{array}{c c c c} f _ {1} (x _ {1}) & f _ {1} (x _ {2}) & \dots & f _ {1} (x _ {n}) \\ f _ {2} (x _ {1}) & f _ {2} (x _ {2}) & \dots & f _ {2} (x _ {n}) \\ \dots & \dots & \dots & \dots \\ f _ {n} (x _ {1}) & f _ {n} (x _ {2}) & \dots & f _ {n} (x _ {n}) \end{array} \right| \vdots \left| \begin{array}{c c c c} g _ {1} (x _ {1}) & g _ {1} (x _ {2}) & \dots & g _ {1} (x _ {n}) \\ g _ {2} (x _ {1}) & g _ {2} (x _ {2}) & \dots & g _ {2} (x _ {n}) \\ \dots & \dots & \dots & \dots \\ g _ {n} (x _ {1}) & g _ {n} (x _ {2}) & \dots & g _ {n} (x _ {n}) \end{array} \right|
\]

is equal to the ratio of the two determinants

\[
\left| \begin{array}{c c c c} f _ {1} (\xi_ {1}) & f _ {1} ^ {\prime} (\xi_ {2}) & \dots & f _ {1} ^ {(n - 1)} (\xi_ {n}) \\ f _ {2} (\xi_ {1}) & f _ {2} ^ {\prime} (\xi_ {2}) & \dots & f _ {2} ^ {(n - 1)} (\xi_ {n}) \\ \dots & \dots & \dots & \dots \\ f _ {n} (\xi_ {1}) & f _ {n} ^ {\prime} (\xi_ {2}) & \dots & f _ {n} ^ {(n - 1)} (\xi_ {n}) \end{array} \right|: \left| \begin{array}{c c c c} g _ {1} (\xi_ {1}) & g _ {1} ^ {\prime} (\xi_ {2}) & \dots & g _ {1} ^ {(n - 1)} (\xi_ {n}) \\ g _ {2} (\xi_ {1}) & g _ {2} ^ {\prime} (\xi_ {2}) & \dots & g _ {2} ^ {(n - 1)} (\xi_ {n}) \\ \dots & \dots & \dots & \dots \\ g _ {n} (\xi_ {1}) & g _ {n} ^ {\prime} (\xi_ {2}) & \dots & g _ {n} ^ {(n - 1)} (\xi_ {n}) \end{array} \right|
\]

where

\[
\xi_ {1} = x _ {1}, \quad \xi_ {1} <   \xi_ {2} <   x _ {2}, \quad \xi_ {2} <   \xi_ {3} <   x _ {3}, \quad \dots , \quad \xi_ {n - 1} <   \xi_ {n} <   x _ {n}
\]

(apply exerc. 9 of I, p.~38).

Particular case where \(g_{1}(x)=1\) , \(g_{2}(x)=x,\ldots,g_{n}(x)=x^{n-1}\) .

¶ 12) \(a)\) Let \(\mathbf{f}\) be a vector function defined and continuous on the finite interval \(\mathrm{I} = [-a, + a]\) , taking its values in a normed space \(\mathrm{E}\) over \(\mathbf{R}\) and twice differentiable on \(\mathrm{I}\) . If one puts \(\mathrm{M}_0 = \sup_{x\in \mathrm{I}}\| \mathbf{f}(x)\|\) , \(\mathrm{M}_2 = \sup_{x\in \mathrm{I}}\| \mathbf{f}''(x)\|\) , show that for all \(x\in \mathrm{I}\) one has

\[
\left\| \mathbf {f} ^ {\prime} (x) \right\| \leqslant \frac {\mathrm{M} _ {0}}{a} + \frac {x ^ {2} + a ^ {2}}{2 a} \mathrm{M} _ {2}
\]

(express each of the differences \(\mathbf{f}(a) - \mathbf{f}(x)\) , \(\mathbf{f}(-a) - \mathbf{f}(x)\) ).

\begin{enumerate}
\def\labelenumi{\alph{enumi})}
\setcounter{enumi}{1}
\tightlist
\item
  Deduce from a) that if f is a twice differentiable function on an interval I (bounded or not), and if \(\mathbf{M}_{0} = \sup_{x \in I} \| \mathbf{f}(x) \|\) and \(\mathbf{M}_{2} = \sup_{x \in I} \| \mathbf{f}''(x) \|\) are finite, then so is \(\mathbf{M}_{1} = \sup_{x \in I} \| \mathbf{f}'(x) \|\) , and one has:
\end{enumerate}

\[
\mathrm{M} _ {1} \leqslant 2 \sqrt {\mathrm{M} _ {0} \mathrm{M} _ {2}} \quad \text { if   I   has   length } \geqslant 2 \sqrt {\frac {M _ {0}}{\mathrm{M} _ {2}}}
\]

\[
\mathrm{M} _ {1} \leqslant \sqrt {2} \sqrt {\mathrm{M} _ {0} \mathrm{M} _ {2}} \quad \text { if } \mathrm{I} = \mathbf {R}.
\]

Show that in these two inequalities the numbers 2 and \(\sqrt{2}\) respectively cannot be replaced by smaller numbers (consider first the case where one supposes merely that f admits a second right derivative, and show that in this case the two terms of the preceding inequalities can become equal, taking for f a real function equal ``in pieces'' to second degree polynomials).

\begin{enumerate}
\def\labelenumi{\alph{enumi})}
\setcounter{enumi}{2}
\tightlist
\item
  Deduce from \(b)\) that if \(\mathbf{f}\) is \(p\) times differentiable on \(\mathbf{R}\) , and if \(\mathbf{M}_p = \sup_{x \in \mathbf{R}} \| \mathbf{f}^{(p)}(x) \|\)
\end{enumerate}

and \(M_{0}=\sup_{x\in\mathbf{R}}\|\mathbf{f}(x)\|\) are finite, then each of the numbers \(M_{k}=\sup_{x\in\mathbf{R}}\left\|\mathbf{f}^{(k)}(x)\right\|\) is finite (for

\(1 \leqslant k \leqslant p - 1)\) and

\[
\mathbf {M} _ {k} \leqslant 2 ^ {k (p - k) / 2} \mathbf {M} _ {0} ^ {1 - k / p} \mathbf {M} _ {p} ^ {k / p}.
\]

¶ 13) a) Let \(f\) be a twice differentiable real function on \(\mathbf{R}\) , such that \((f(x))^2 \leqslant a\) and \((f'(x))^2 + (f''(x))^2 \leqslant b\) on \(\mathbf{R}\) ; show that

\[
(f (x)) ^ {2} + \left(f ^ {\prime} (x)\right) ^ {2} \leqslant \max (a, b)
\]

on R (argue by contradiction, noting that if the function \(f^{2}+f^{\prime2}\) takes a value \(c > \max(a, b)\) at a point \(x_{0}\) then there exist two points \(x_{1}, x_{2}\) such that \(x_{1} < x_{0} < x_{2}\) and that at \(x_{1}\) and \(x_{2}\) the function \(f^{\prime}\) takes values small enough that \(f^{2} + f^{\prime2}\) takes values \textless{} c; then consider a point of \([x_{1}, x_{2}]\) where \(f^{2} + f^{\prime2}\) attains its supremum on this interval).

\begin{enumerate}
\def\labelenumi{\alph{enumi})}
\setcounter{enumi}{1}
\tightlist
\item
  Let \(f\) be a real function \(n\) times differentiable on \(\mathbf{R}\) , and such that \((f(x))^2 \leqslant a\) and \((f^{(n - 1)}(x))^2 + (f^{(n)}(x))^2 \leqslant b\) on \(\mathbf{R}\) ; show that then
\end{enumerate}

\[
(f ^ {(k - 1} (x)) ^ {2} + (f ^ {(k)} (x)) ^ {2} \leqslant \max (a, b)
\]

on \(\mathbf{R}\) for \(1 \leqslant k \leqslant n\) . (Argue by induction on \(n\) ; note that, by exerc. 12 the supremum \(c\) of \((f'(x))^2\) on \(\mathbf{R}\) is finite; show that one necessarily has \(c \leqslant \max(a, b)\) by reducing to a contradiction: assuming that \(c > \max(a, b)\) choose the constants \(\lambda\) and \(\mu\) so that for the function \(g = \lambda f + \mu\) one has \(|g(x)| \leqslant 1\) , \(|g'(x)| \leqslant 1\) , yet one cannot have \((g(x))^2 + (g'(x))^2 \leqslant 1\) for all \(x\) .

¶ 14) Let \(\mathbf{f}\) be a function which is \(n - 1\) times differentiable on an interval I containing 0, and let \(\mathbf{f}_n\) be the vector function defined for \(x \neq 0\) on I by the relation

\[
\mathbf {f} (x) = \mathbf {f} (0) + \mathbf {f} ^ {\prime} (0) \frac {x}{1 !} + \mathbf {f} ^ {\prime \prime} (0) \frac {x ^ {2}}{2 !} + \dots + \mathbf {f} ^ {(n - 1)} (0) \frac {x ^ {n - 1}}{(n - 1) !} + \mathbf {f} _ {n} (x) x ^ {n}.
\]

\begin{enumerate}
\def\labelenumi{\alph{enumi})}
\item
  Show that if \(\mathbf{f}\) has an \((n + p)^{th}\) derivative at the point 0 then \(\mathbf{f}_n\) has a \(p^{th}\) derivative at the point 0 and an \((n + p - 1)^{th}\) derivative at all points of a neighbourhood of 0 distinct from 0; moreover, one has \(\mathbf{f}_n^{(k)}(0) = \frac{k!}{(n + k)!}\mathbf{f}^{(n + k)}(0)\) for \(0 \leqslant k \leqslant p\) , and \(\mathbf{f}_n^{(p + k)}(x)x^k\) tends to 0 with \(x\) , for \(1 \leqslant k \leqslant n - 1\) (express the derivatives of \(\mathbf{f}_n\) with the help of the Taylor expansions of the successive derivatives of \(\mathbf{f}\) , and use prop. 6 of I, p.~18).
\item
  Conversely, let \(f_{n}\) be a vector function admitting an \((n + p - 1)^{th}\) derivative on a neighbourhood of 0 in I, and such that \(\mathbf{f}_{n}^{(p+k)}(x)x^{k}\) has a limit for \(0 \leqslant k \leqslant n - 1\) . Show that the function \(\mathbf{f}_{n}(x)x^{n}\) has an \((n + p - 1)^{th}\) derivative on I; if, further, \(f_{n}\) admits a \(p^{th}\) derivative at the point 0, then \(\mathbf{f}_{n}(x)x^{n}\) admits an \((n + p)^{th}\) derivative at the point 0.
\item
  Suppose that I is symmetric with respect to 0 and that f is even ( \(\mathbf{f}(-x) = \mathbf{f}(x)\) on I). Show, with the help of a), that if f is 2n times differentiable on I, then there exists a function g defined and n times differentiable on I, such that \(\mathbf{f}(x) = \mathbf{g}(x^{2})\) on I.
\end{enumerate}

¶ 15) Let I be an open interval in \(\mathbf{R}\) , and \(\mathbf{f}\) a vector function defined and continuous on I; suppose that there are \(n\) vector functions \(\mathbf{g}_i\) ( \(1 \leqslant i \leqslant n\) ) defined on I, and such that the function of \(x\)

\[
\frac {1}{h ^ {n}} \left(\mathbf {f} (x + h) - \mathbf {f} (x) - \sum_ {p = 1} ^ {n} \frac {h ^ {p}}{p !} \mathbf {g} _ {p} (x)\right)
\]

tends uniformly to 0 on every compact interval contained in I as h tends to 0.

\begin{enumerate}
\def\labelenumi{\alph{enumi})}
\item
  We put \(\mathbf{f}_p(x,h) = \Delta^p\mathbf{f}(x;h,h,\ldots ,h)\) (I, p.~40, exerc. 6). Show that, for \(1\leqslant p\leqslant n\) , \((1 / h^{p})\mathbf{f}_{p}(x,h)\) tends uniformly to \(\mathbf{g}_p(x)\) on every compact subinterval of I as \(h\) tends to 0, and that the \(\mathbf{g}_p\) are continuous on I (prove this successively for \(p = n\) , \(p = n - 1\) , etc.)
\item
  Deduce from this that \(\mathbf{f}\) has a continuous \(n^{th}\) derivative and that \(\mathbf{f}^{(p)} = \mathbf{g}_p\) for \(1 \leqslant p \leqslant n\) (taking account of the relation \(\mathbf{f}_{p+1}(x, h) = \mathbf{f}(x + h, h) - \mathbf{f}_p(x, h)\) ).
\end{enumerate}

¶ 16) Let \(f\) be a real function \(n\) times differentiable on \(\mathrm{I} = ] - 1, + 1[\) , and such that \(|f(x)|\leqslant 1\) on this interval.

\begin{enumerate}
\def\labelenumi{\alph{enumi})}
\tightlist
\item
  Show that if \(m_k(\lambda)\) denotes the minimum of \(|f^{(k)}(x)|\) on an interval of length \(\lambda\) contained in I then one has
\end{enumerate}

\[
m _ {k} (\lambda) \leqslant \frac {2 ^ {k (k + 1) / 2} k ^ {k}}{\lambda^ {k}}
\]

\[
(1 \leqslant k \leqslant n).
\]

(Note that if the interval of length \(\lambda\) is decomposed into three intervals of lengths \(\alpha, \beta, \gamma\) , one has

\[
m _ {k} (\lambda) \leqslant \frac {1}{\beta} (m _ {k - 1} (\alpha) + m _ {k - 1} (\gamma)).
\]

\begin{enumerate}
\def\labelenumi{\alph{enumi})}
\setcounter{enumi}{1}
\tightlist
\item
  Deduce from a) that there exists a number \(\mu_{n}\) depending only on the integer n such that if \(|f'(0)| \geqslant \mu_{n}\) , then the derivative \(f^{(n)}(x)\) vanishes on at least n - 1 distinct points of I (show by induction on k that \(f^{(k)}\) vanishes at least k - 1 times on I).
\end{enumerate}

\begin{enumerate}
\def\labelenumi{\arabic{enumi})}
\setcounter{enumi}{16}
\tightlist
\item
  \begin{enumerate}
  \def\labelenumii{\alph{enumii})}
  \tightlist
  \item
    Let f be a vector function having derivatives of all orders on an open interval \(I \subset R\) . Suppose that, on I, one has \(\left\|\mathbf{f}^{(n)}(x)\right\| \leqslant a n!r^n\) , where a and r are two numbers \textgreater0 and independent of x and n; show that at each point \(x_0\) the ``Taylor series'' with general term \((1/n!) \mathbf{f}^{(n)}(x_0)(x - x_0)^n\) is convergent, and has sum \(\mathbf{f}(x)\) on some neighbourhood of \(x_0\) .
  \end{enumerate}
\end{enumerate}

\begin{enumerate}
\def\labelenumi{\alph{enumi})}
\setcounter{enumi}{1}
\item
  Conversely, if the Taylor series for \(\mathbf{f}\) at a point \(x_0\) converges on a neighbourhood of \(x_0\) there exist two numbers \(a\) and \(r\) (depending on \(x_0\) ) such that \(\left\| \mathbf{f}^{(n)}(x_0)\right\| \leqslant a.n!r^n\) for every integer \(n > 0\) .
\item
  Deduce from a) and exerc. 16 b) that if, on an open interval \(\mathrm{I} \subset \mathbf{R}\) , a real function \(f\) is indefinitely differentiable and if there is an integer \(p\) independent of \(n\) such that, for all \(n\) , the function \(f^{(n)}\) does not vanish at more than \(p\) distinct points of \(\mathrm{I}\) , then the Taylor series of \(f\) on a neighbourhood of each point \(x_0 \in \mathrm{I}\) is convergent, and has sum \(f(x)\) at every point of a neighbourhood of \(x_0\) .
\end{enumerate}

\begin{enumerate}
\def\labelenumi{\arabic{enumi})}
\setcounter{enumi}{17}
\tightlist
\item
  Let \((a_{n})_{n\geqslant 0}\) be an arbitrary sequence of complex numbers. For each \(n\geqslant 0\) put \(s_n^{(0)} = a_n\) , and, inductively, for \(k\geqslant 0\) , define
\end{enumerate}

\[
s _ {n} ^ {(k + 1)} = s _ {0} ^ {(k)} + s _ {1} ^ {(k)} + \dots + s _ {n - 1} ^ {(k)}.
\]

\begin{enumerate}
\def\labelenumi{\alph{enumi})}
\tightlist
\item
  Prove ``Taylor's formula for sequences'': for each integer
\end{enumerate}

\[
\left| s _ {n + h} ^ {(k)} - s _ {n} ^ {(k)} - h s _ {n} ^ {(k - 1)} - \binom {h} {2} s _ {n} ^ {(k - 2)} - \dots - \binom {h} {k - 1} s _ {n} ^ {(1)} \right| \leqslant \binom {h} {k} \sup _ {0 \leqslant j \leqslant h - 1} | a _ {n + j} |
\]

(proceed by induction on k).

\begin{enumerate}
\def\labelenumi{\alph{enumi})}
\setcounter{enumi}{1}
\tightlist
\item
  Suppose that there is a number C such that \(|na_{n}| \leqslant C\) for all n, and that the sequence \((s^{(2)}/n)\) formed by the arithmetic means \((s_{0} + \cdots + s_{n-1})/n\) of the partial sums \(s_{n} = a_{0} + \cdots + a_{n-1}\) tends to a limit \(\sigma\) . Show that the series with general term \(a_{n}\) is convergent and has sum \(\sigma\) (``Hardy-Littlewood tauberian theorem''). (Write
\end{enumerate}

\[
s _ {n} = \frac {1}{h} \left(s _ {n + h} ^ {(2)} - s ^ {(2)}\right) + \frac {h - 1}{2} r _ {n}
\]

where \(|r_n|\) is bounded above with the aid of the inequality \(|na_n| \leqslant C\) , and \(h\) is chosen suitably as a function of \(n\) .)

\setcounter{exercisegroup}{3}
\refstepcounter{exercisegroup}
